%% Dans ce fichier, les �ventuels accents sont cod�s
%% en ISO-8859-15 (aka latin9)
%% La ligne suivante donne trois exemples 
%% e accent aigu : �, e accent grave : �, e-dans-l'o : �
%%%%%ENTREES%%%%
\begin{tabarg}
\multicolumn{3}{l}{\titretab {Entr\'ees}}\\
\hline

$X$ & Vecteur, \texttt{float[$N_{level}$]} & Signal lidar \\
$N_{level}$ & Coeff., \texttt{float} & Nombre de portes (\$Arg1.array\_dim)\\
$X_{gate}$ & Coeff., \texttt{String} & Bornes ou portes du profil (\$Arg1.bornes) (m)\\
$Z$ & Vecteur, \texttt{float[$N_{level}$]} & Altitudes r�elles des portes du lidar (m)\\
$\beta_{mol}$ & Vecteur, \texttt{float[$N_{level}$]} &  Profil du coefficient de r�tro-diffusion mol�culaire ($m^{-1}$)\\
$Alt_{nuage}$ & Vecteur, \texttt{float} & Altitude des nuages (vaut 0 si pas de nuage) (m)\\
$AOD$ & Vecteur, \texttt{float} & Epaisseur optique des a�rosols\\
$Z_{ref}$ & Vecteur, \texttt{float} & Altitude de r�f�rence pour le signal mol�culaire (m)\\
$D_{min}$ & Vecteur, \texttt{float} & Distance minimale du lidar � partir de laquelle on d�termine le profil (m)\\
$\beta_{ainit}$ & Coeff., \texttt{float} & Valeur d'initialisation de $\beta$  ($m^{-1}$)\\
$Extrap$ & Coeff., \texttt{String} & Type d'extrapolation pour la zone proche du lidar (< $D_{min}$) :"lin" ou "cst" \\
\end{tabarg}

\medskip
%%%%%SORTIES%%%%
\begin{tabarg}
\multicolumn{3}{l}{\titretab {Sorties}}\\
\hline
$\sigma$ & Vecteur, \texttt{float[$N_{level}$]} & Coefficient d'extinction des a�rosols (m-1)\\
$\beta$ & Vecteur, \texttt{float[$N_{level}$]} & Coefficient de r�tro-diffusion des a�rosols (m-1)\\
$SA$ & Vecteur, \texttt{float} & Rapport lidar\\
\end{tabarg}
